\chapter{Introduction} \label{ch:introduction}

\section{Motivation and Background}
The ongoing shift from fossil fuels to renewable energy sources (RES) brings new challenges with it. Weather dependent RES like solar or wind energy have high generation variability, which has a negative impact on the power system, creating reliability on the controllable, traditional power generation by fossils. To counteract the negatives, stabilize the generation output and therefore complete the energy transition, stakeholders are more and more implementing energy storage systems (ESS) projects into the power grid. One of the fastest acting energy storage technologies (EST) with high volumetric energy density, high power density and round-trip efficiency are battery energy storage systems (BESS) like Lithium-ion batteries. For these reasons and due to favorable price-development, BESS make the second biggest operational EST capacity globally, following pumped hydro \cite{odero_comprehensive_2022}. With fast growing BESS capacity, there is high potential for optimization in terms of sizing and siting future BESS projects in order to reach higher technical and economic efficiency and finally speed up the energy transition.

\section{Core Objectives and Research Question}
This thesis aims to develop an optimization-based modeling framework to determine optimal allocation strategies for Battery Energy Storage Systems, considering both technical and economic constraints for transmission grid operators.

The core objective is to determine the optimal generation and storage capacity of BESS -- expressed in MW and MWh, respectively -- at the nodal grid level.

To guide this analysis and bridge the identified gaps in current energy system modeling, the following exploratory research question is posed: 
\begin{quote}
\textit{What insights into strategic investor behavior can be gained by shifting from a traditional central-planner cost-minimization approach to a decentralized, game-theoretic EPEC framework when determining optimal BESS capacities?}
\end{quote}

\section{Method Applied}
To fulfill these objectives, this thesis departs from traditional central-planner cost-minimization approaches and adopts a game-theoretic perspective. In liberalized electricity markets, investment decisions are made by independent, profit-maximizing actors rather than a single entity, as highlighted by Conejo and Ruiz \cite{conejo_complementarity_2020}. Therefore, the interactions of strategic BESS investors are modeled as an Equilibrium Problem with Equilibrium Constraints (EPEC).

An EPEC is a set of interrelated bi-level optimization problems, known as Mathematical Programs with Equilibrium Constraints (MPECs) \cite{conejo_complementarity_2020}. In this framework, the upper level represents the strategic BESS investors maximizing their individual Net Present Value (NPV), while the lower level represents the Independent System Operator (ISO) clearing the market via a cost-minimal direct current optimal power flow (DC-OPF). Similar EPEC formulations have been successfully applied in the energy system context to analyze equilibria in coupled electricity and natural gas markets, demonstrating the method's capability to capture complex, strategic market dynamics \cite{chen_equilibria_2020}.

To calculate the cost-minimal and profit-maximizing investment decisions, the proposed mathematical optimization model is implemented in Python. The complex mathematical formulation of the non-linear complementarity conditions is built using the Pyomo optimization package \cite{pyomo_2021}. The resulting non-linear programming (NLP) problems are subsequently solved using the state-of-the-art interior-point optimizer Ipopt \cite{ipopt_2006}, utilizing a diagonalization algorithm to compute the Nash Equilibria among the investors.

\section{Outline of the Master Thesis}
This work is organized as follows. Chapter \ref{ch:state_of_the_art} provides a comprehensive state-of-the-art review of BESS allocation in power grids and the application of EPECs in energy system modeling, defining the novelty and progress beyond the current literature. Chapter \ref{ch:Methodology} elaborates on the applied EPEC methodology, detailing the mathematical formulation, the model expansions and the algorithmic solution strategy. Chapter \ref{ch:results} presents the quantitative results and discusses the strategic investment behavior, grid congestion, and the "Negative Arbitrage" phenomenon. Chapter \ref{ch:synthesis} synthesizes these findings, answers the core research question, and critically evaluates the strengths and limitations of the EPEC approach. Finally, Chapter \ref{ch:conclusion} concludes the thesis and provides an outlook for future research.


